%% 2i2d-tour-guet — élévation de la tour de guet (exercice guidé) : pylône en
%% treillis sur deux massifs A et B, cabine et panneaux photovoltaïques en tête,
%% poids P au centre de gravité G et vent F appliqué en V.
%% Échelle verticale comprimée (double barre de rupture sur le fût).
\documentclass[tikz,border=4pt]{standalone}
\usepackage{liaisons}
\begin{document}
\begin{tikzpicture}[x=1cm,y=0.85cm,   % échelle verticale comprimée
  force/.style={-{Stealth[length=5pt]}, line width=1pt, solideA},
  barre/.style={line width=0.5pt, solideB},
  cote/.style={{Stealth[length=4pt]}-{Stealth[length=4pt]}, line width=0.6pt, solideD},
  attache/.style={line width=0.4pt, solideD!80},
  etiq/.style={font=\small, inner sep=2pt},
  petit/.style={font=\scriptsize, inner sep=1.5pt}]

\newcommand\batihoriz[3]{%
  \draw[line width=1pt, solideE] (#1,#3) -- (#2,#3);
  \begin{scope}
    \clip (#1,#3-0.26) rectangle (#2,#3);
    \foreach \hx in {-16,-15.76,...,16} { \draw[line width=0.5pt, solideE!75] (\hx,#3) -- (\hx-0.3,#3-0.3); }
  \end{scope}}

%% --- abscisses des montants en fonction de la hauteur (le pylône se rétrécit) ---
\newcommand\mg[1]{2.4+0.225*#1}   % montant gauche
\newcommand\md[1]{5.2-0.225*#1}   % montant droit

%% --- sol et massifs de fondation ------------------------------------------------
\batihoriz{0.5}{7.5}{0}
\draw[line width=0.7pt, draw=solideE, fill=solideE!20] (2.15,-0.275) rectangle (2.65,-0.025);
\draw[line width=0.7pt, draw=solideE, fill=solideE!20] (4.95,-0.275) rectangle (5.45,-0.025);
\node[etiq, anchor=north] at (2.4,-0.3) {$B$};
\node[etiq, anchor=north] at (5.2,-0.3) {$A$};

%% --- fût en treillis --------------------------------------------------------------
\draw[line width=1pt, solideB] (2.4,0) -- (3.3,4.0);
\draw[line width=1pt, solideB] (5.2,0) -- (4.3,4.0);
\foreach \y in {0.8,1.6,2.4,3.2} { \draw[barre] (\mg{\y},\y) -- (\md{\y},\y); }   % traverses
\draw[barre] (\mg{0},0) -- (\md{0.8},0.8);        % diagonales en zigzag
\draw[barre] (\md{0.8},0.8) -- (\mg{1.6},1.6);
\draw[barre] (\mg{1.6},1.6) -- (\md{2.4},2.4);
\draw[barre] (\md{2.4},2.4) -- (\mg{3.2},3.2);
\draw[barre] (\mg{3.2},3.2) -- (\md{4.0},4.0);
\draw[barre] (\mg{4.0},4.0) -- (\md{4.0},4.0);

%% --- double barre de rupture (échelle verticale comprimée) -------------------------
\draw[line width=0.9pt, solideD] (2.98,2.52) -- (4.80,2.74);
\draw[line width=0.9pt, solideD] (2.94,2.74) -- (4.76,2.96);
\node[petit, anchor=east, text=solideD, align=right] at (2.86,2.7) {échelle\\comprimée};

%% --- cabine, toiture et panneaux ------------------------------------------------------
\draw[line width=0.8pt, draw=solideB, fill=solideB!10] (3.0,4.0) rectangle (4.6,4.55);
\draw[line width=0.8pt, draw=solideB, fill=solideB!18, line join=round]
  (2.85,4.55) -- (4.75,4.55) -- (4.25,4.95) -- (3.35,4.95) -- cycle;
\draw[line width=0.8pt, draw=solideF, fill=solideF!30] (3.0,3.85) rectangle (4.6,4.0);

\draw[attache] (4.62,3.92) -- (5.25,3.25);
\node[petit, anchor=west, text=solideF, align=left] at (5.28,3.25) {14 panneaux\\10 m$^2$};
\draw[attache] (4.68,4.66) -- (5.25,4.72);
\node[petit, anchor=west, text=solideB, align=left] at (5.28,4.72) {cabine + toiture\\4 m$^2$};

%% --- le vent appliqué en V --------------------------------------------------------------
\fill (3.0,4.0) circle (1.5pt);
\node[etiq, anchor=north east, inner sep=1pt] at (2.96,3.93) {$V$};
\draw[force] (1.8,4.0) -- (2.96,4.0);
\node[etiq, anchor=south, text=solideA, align=center] at (2.25,4.08) {$\vec{F}$ (vent)};

%% --- le poids appliqué en G ---------------------------------------------------------------
\draw[line width=0.6pt, solideF, dash pattern=on 6pt off 2pt on 1pt off 2pt] (3.8,-0.4) -- (3.8,4.85);
\fill (3.8,2.0) circle (1.5pt);
\node[etiq, anchor=east, fill=white, inner sep=1pt] at (3.72,2.02) {$G$};
\draw[force] (3.8,2.0) -- (3.8,1.1);
\node[etiq, anchor=west, text=solideA, fill=white, inner sep=1.5pt] at (3.88,1.45) {$\vec{P} = 55$ kN};

%% --- niveaux et cotes -------------------------------------------------------------------------
\draw[attache] (4.78,4.0) -- (5.25,4.0);
\node[petit, anchor=west, text=solideD] at (5.28,4.0) {+ 40,70};
\draw[attache] (5.5,0) -- (6.55,0);
\node[petit, anchor=south east, text=solideD] at (6.55,0.04) {$\pm$ 0,00};

\draw[attache] (3.8,-0.45) -- (3.8,-0.8)  (5.2,-0.45) -- (5.2,-1.2)  (2.4,-0.45) -- (2.4,-1.2);
\draw[cote] (3.8,-0.66) -- (5.2,-0.66);
\node[petit, anchor=south, text=solideD] at (4.5,-0.64) {4 m};
\draw[cote] (2.4,-1.08) -- (5.2,-1.08);
\node[petit, anchor=south, text=solideD] at (3.8,-1.06) {8 m};

\node[petit, anchor=north, text=black!70] at (4.0,-1.32)
     {à la limite du basculement, l'action en $B$ est nulle};
\end{tikzpicture}
\end{document}
